% =====================================================================
%  Chapter 2: QCD on the Lattice - a First Look
%  Gattringer & Lang, "Quantum Chromodynamics on the Lattice" (LNP 788)
%  Reconstructed from extract/ch02.txt; equations auto-numbered.
% =====================================================================

\chapter{QCD on the Lattice -- a First Look}

In Chap.~1 we derived the lattice path integral for a scalar field theory
using the canonical approach for the quantization of the system.
Subsequently, we changed our point of view and adopted the lattice path
integral itself as a method of quantization. The steps involved are a
lattice discretization of the classical Euclidean action and the
construction of the measure for the integration over ``all configurations''
of the classical fields. In this chapter we begin with the implementation of
these two steps for quantum chromodynamics.

Quantum chromodynamics (QCD) is the theory of strongly interacting particles
and fields, i.e., the theory of quarks and gluons. In the first section of
this chapter we will review the QCD action functional and its symmetries.
This serves as a preparation for the subsequent discretization of QCD on the
lattice. The discretization proceeds in several steps. We begin with the
naive discretization of the fermionic part of the QCD action followed by
discussing the lattice action for the gluons. In the end we write down a
complete expression for the QCD lattice path integral. Already at this point
we stress, however, that the expression for the path integral obtained in
this chapter is not the final word. The nature of the quark fields in our
path integral will have to be changed in order to incorporate Fermi
statistics. This step will be taken in Chap.~5 where we reinterpret the
quark fields as anti-commuting numbers, so-called Grassmann numbers.

\section{The QCD action in the continuum}

We begin our construction of the QCD lattice path integral with a review of
the Euclidean continuum action. After introducing the quark and gluon
fields, we develop the QCD continuum action starting with the discussion of
the fermionic part of the action and its invariance under gauge
transformations. Subsequently, we construct an action for the gluon fields
that are invariant under these transformations.

\subsection{Quark and gluon fields}

The quarks are massive fermions and as such are described by Dirac 4-spinors
\begin{equation}
  \psi^{(f)}(x)^\alpha_c , \qquad \bar\psi^{(f)}(x)^\alpha_c .
\end{equation}
These quark fields carry several indices and arguments. The space--time
position is denoted by $x$, the Dirac index by $\alpha = 1, 2, 3, 4$, and
the color index by $c = 1, 2, 3$. In general we will use Greek letters for
Dirac indices and letters $a, b, c, \ldots$ for color. Each field
$\psi^{(f)}(x)$ thus has 12 independent components. In addition the quarks
come in several flavors called up, down, strange, charm, bottom, and top,
which we indicate by a flavor index $f = 1, 2, \ldots, 6$. In many
calculations it is sufficient to include only the lightest two or three
flavors of quarks. Thus, our flavor index will run from 1 to $\Nf$, the
number of flavors. We remark that often we omit the indices and use
matrix/vector notation instead.

In the Minkowskian operator approach the fields $\psi$ and $\bar\psi$ are
related by $\bar\psi = \psi^\dagger \gamma_0$, where $\gamma_0$ is the
$\gamma$-matrix related to time. In the Euclidean path integral one uses
independent integration variables $\psi$ and $\bar\psi$.

In addition to the quarks, QCD contains gauge fields describing the gluons,
\begin{equation}
  A_\mu(x)_{cd} .
\end{equation}
These fields also carry several indices. As for the quark fields we have a
space--time argument denoted by $x$. In addition, the gauge fields
constitute a vector field carrying a Lorentz index $\mu$ which labels the
direction of the different components in space--time. Since we are
interested in the Euclidean action, the Lorentz index $\mu$ is Euclidean,
i.e., we do not distinguish between covariant and contravariant indices.
There is no metric tensor involved and $\mu = 1, 2, 3, 4$ simply label the
different components. Finally, the gluon field carries color indices
$c, d = 1, 2, 3$. For given $x$ and $\mu$, the field $A_\mu(x)$ is a
traceless, hermitian $3 \times 3$ matrix at each space--time point $x$. We
will discuss the structure of these matrices and their physical motivation
in more detail below.

It is convenient to split the QCD action into a fermionic part, which
includes quark fields and an interaction term coupling them to the gluons,
and a gluonic part, which describes propagation and interaction of only the
gluons.

\subsection{The fermionic part of the QCD action}

The fermionic part $S_F[\psi, \bar\psi, A]$ of the QCD action is a bilinear
functional in the fields $\psi$ and $\bar\psi$. It is given by
\begin{equation}
  \begin{split}
    S_F[\psi, \bar\psi, A] &=
      \sum_{f=1}^{\Nf} \int \dd^4 x\,
      \bar\psi^{(f)}(x) \Big(\gamma_\mu \big(\partial_\mu
        + \im A_\mu(x)\big) + m^{(f)}\Big) \psi^{(f)}(x) \\
    &= \sum_{f=1}^{\Nf} \int \dd^4 x\,
      \bar\psi^{(f)}(x)^\alpha_c \Big( (\gamma_\mu)_{\alpha\beta}
      \big(\delta_{cd}\,\partial_\mu + \im A_\mu(x)_{cd}\big)\\
    &\qquad + m^{(f)}\,\delta_{\alpha\beta}\,\delta_{cd}\Big)
      \psi^{(f)}(x)^\beta_d .
  \end{split}
\end{equation}
In the first line of this equation we have used matrix/vector notation for
the color and Dirac indices, while in the second line we write all indices
explicitly. Note that we use the Einstein summation convention.

Equation (2.3) makes it obvious that the action is a sum of the actions for
the individual flavors $f = 1, 2, \ldots, \Nf$. The quarks with different
flavor all couple in exactly the same way to the gluon field $A_\mu$ and
only differ in their masses $m^{(f)}$. Of course, different flavors also
have different electric charge and thus couple differently to the
electromagnetic field. However, here we only discuss the strong interaction.

The color indices $c$ and $d$ of the quark fields $\psi$, $\bar\psi$ are
summed over with the corresponding indices of the gauge field and, in this
way, couple the quarks to the gluons. The coupling of the gauge field is
different for each component $\mu$ since each component is multiplied with
a different matrix $\gamma_\mu$. The $\gamma$-matrices are $4\times 4$
matrices in Dirac space, and in the QCD action they mix the different Dirac
components of the quark fields. They are the Euclidean versions of the
(Minkowski) $\gamma$-matrices familiar from the Dirac equation. The
Euclidean $\gamma$-matrices $\gamma_\mu$, $\mu = 1, 2, 3, 4$, obey the
Euclidean anti-commutation relations
\begin{equation}
  \{\gamma_\mu, \gamma_\nu\} = 2\,\delta_{\mu\nu}\,\mathbb{1} .
\end{equation}
In Appendix A.2 we discuss how to construct the Euclidean $\gamma$-matrices
from their Minkowski counterparts and give an explicit representation. The
different partial derivatives in the kinetic term of (2.3) mix the Dirac
components in the same way as the gauge fields, i.e., the $\partial_\mu$ are
also contracted with the matrices $\gamma_\mu$. The kinetic term is,
however, trivial in color space. The mass term, finally, is trivial in both
color and Dirac space.

Having discussed our notation in detail, we still should verify that the
action (2.3) indeed gives rise to the relativistic wave equation for
fermions, the Dirac equation. For a single flavor, the contribution to the
action is given by (we drop the flavor index for the subsequent discussion
and use matrix/vector notation for the color and Dirac indices)
\begin{equation}
  S_F[\psi, \bar\psi, A] = \int \dd^4 x\,
    \bar\psi(x) \big(\gamma_\mu (\partial_\mu + \im A_\mu(x)) + m\big)
    \psi(x) .
\end{equation}
The simplest way of applying the Euler--Lagrange equations (1.43) in (2.5)
is to differentiate the integrand of (2.5) with respect to $\bar\psi(x)$.
This gives rise to
\begin{equation}
  \big(\gamma_\mu (\partial_\mu + \im A_\mu(x)) + m\big) \psi(x) = 0 ,
\end{equation}
which is indeed the (Euclidean) Dirac equation in an external field
$A_\mu$. Thus, we have verified that the action (2.5) has the correct form.

\subsection{Gauge invariance of the fermion action}

So far we have only discussed the different building blocks of QCD and how
they are assembled in the fermionic part of the QCD action. Let us now dive
a little bit deeper into the underlying structures and symmetries.

Up to the additional color structure, the action (2.5) is exactly the action
of electrodynamics -- when using matrix/vector notation this difference is
not even explicit. As a matter of fact, the QCD action is obtained by
generalizing the gauge invariance of electrodynamics.

In electrodynamics the action is invariant under multiplication of the
fermion fields with an arbitrary phase at each space--time point $x$,
combined with a transformation of the gauge field. In QCD we require
invariance under local rotations among the color indices of the quarks. At
each space--time point $x$ we choose an independent complex $3 \times 3$
matrix $\Omega(x)$. The matrices are required to be unitary,
$\Omega(x)^\dagger = \Omega(x)^{-1}$, and to have $\det[\Omega(x)] = 1$.
Such matrices are the defining representation of the special unitary group,
denoted by SU(3) for the case of $3 \times 3$ matrices. It is easy to see
that this set is closed under matrix multiplication. Furthermore the unit
matrix is contained in this set and for every element there exists an
inverse (the hermitian conjugate matrix). Thus the set of SU(3) matrices
forms a group. We collect the basic equations showing these statements in
Appendix A.1. Note, however, that the group operation -- the matrix
multiplication -- is not commutative. Groups with a non-commutative group
operation are called non-abelian groups. The idea of using non-abelian
groups for a gauge theory was pursued by Yang and Mills [1], and such
theories are often referred to as Yang--Mills theories.

Returning to our discussion of the QCD gauge invariance, we require that the
action is invariant under the transformation
\begin{equation}
  \psi(x) \to \psi'(x) = \Omega(x)\psi(x) , \qquad
  \bar\psi(x) \to \bar\psi'(x) = \bar\psi(x)\Omega(x)^\dagger
\end{equation}
for the fermion fields and a yet unspecified transformation
$A_\mu(x) \to A'_\mu(x)$ for the gauge fields. Invariance of the action
means that we require
\begin{equation}
  S_F[\psi', \bar\psi', A'] = S_F[\psi, \bar\psi, A] .
\end{equation}
With (2.5) and (2.7), this gives
\begin{equation}
  S_F[\psi', \bar\psi', A'] =
    \int \dd^4 x\,
    \bar\psi(x)\Omega(x)^\dagger
    \Big(\gamma_\mu \big(\partial_\mu + \im A'_\mu(x)\big) + m\Big)
    \Omega(x)\psi(x) .
\end{equation}
Using $\Omega(x)^\dagger = \Omega(x)^{-1}$, we see immediately that for the
mass term the gauge transformation matrices cancel. For the other terms the
situation is a bit more involved. From comparing (2.5) with (2.9) we obtain
the condition
\begin{equation}
  \begin{split}
    \partial_\mu + \im A_\mu(x) &=
      \Omega(x)^\dagger \big(\partial_\mu + \im A'_\mu(x)\big) \Omega(x) \\
    &= \partial_\mu + \Omega(x)^\dagger \big(\partial_\mu \Omega(x)\big)
      + \im\, \Omega(x)^\dagger A'_\mu(x)\Omega(x) .
  \end{split}
\end{equation}
This is an equation for an operator acting on a function of $x$. Thus, due
to the product rule, we find two terms with derivatives. We now can solve
for $A'_\mu(x)$ (again we use $\Omega(x)^\dagger = \Omega(x)^{-1}$) and we
arrive at the transformation property for the gauge field
\begin{equation}
  A_\mu(x) \to A'_\mu(x)
  = \Omega(x) A_\mu(x) \Omega(x)^\dagger
    + \im\, \big(\partial_\mu \Omega(x)\big) \Omega(x)^\dagger .
\end{equation}
Note that also $A'_\mu(x)$ is hermitian and traceless as required for the
gauge fields. For the first term on the right-hand side of (2.11), this
follows from the fact that $A_\mu(x)$ is traceless and
$\Omega(x)^\dagger = \Omega(x)^{-1}$. For the second term this is shown in
Appendix A.1 (see (A.11) and (A.15)).

The requirement that the fermion action (2.5) remains invariant under the
gauge transformation (2.7) of the fermions necessarily implies the presence
of gauge fields $A_\mu(x)$ with the transformation properties given by
(2.11).

In the next section, when we discretize QCD on the lattice, we again
require the invariance of the lattice action under the local gauge
transformations (2.7) for the quark fields. Gauge fields have to be
introduced to achieve gauge invariance of the action.

\subsection{The gluon action}

Let us now discuss the action for the gluon fields $A_\mu(x)$. The gluon
action $S_G[A_\mu]$ is a functional of only the gauge fields and is required
to be invariant under the transformation (2.11):
\begin{equation}
  S_G[A'] = S_G[A] .
\end{equation}
To construct an action with this property we define the covariant derivative
\begin{equation}
  D_\mu(x) = \partial_\mu + \im A_\mu(x) .
\end{equation}
From our intermediate result in the first line of (2.10) we read off the
transformation property for the covariant derivative as
\begin{equation}
  D_\mu(x) \to D'_\mu(x) = \partial_\mu + \im A'_\mu(x)
  = \Omega(x) D_\mu(x) \Omega(x)^\dagger .
\end{equation}
These transformation properties ensure that $D_\mu(x)\psi(x)$ and
$\psi(x)$ transform in exactly the same way -- thus, the name ``covariant
derivative.''

The covariant derivatives are now used to construct an action functional
which is a generalization of the expression known from electrodynamics. We
define the field strength tensor $F_{\mu\nu}(x)$ as the commutator
\begin{equation}
  F_{\mu\nu}(x) = -\im [D_\mu(x), D_\nu(x)]
  = \partial_\mu A_\nu(x) - \partial_\nu A_\mu(x)
    + \im [A_\mu(x), A_\nu(x)] ,
\end{equation}
where the last term on the right-hand side does not vanish since $A_\mu(x)$
and $A_\nu(x)$ are matrices, i.e., objects where multiplication is a
non-commutative operation. Up to this commutator, the field strength tensor
$F_{\mu\nu}(x)$ has the same form as the field strength in electrodynamics.

The fact that the field strength tensor is the commutator of two covariant
derivatives implies that it inherits the transformation properties (2.14),
i.e., it transforms as
\begin{equation}
  F_{\mu\nu}(x) \to F'_{\mu\nu}(x) = \Omega(x) F_{\mu\nu}(x) \Omega(x)^\dagger .
\end{equation}
As a candidate for the gauge action we now consider the expression
\begin{equation}
  S_G[A] = \frac{1}{2 g^2} \int \dd^4 x\,
    \tr\big[F_{\mu\nu}(x) F_{\mu\nu}(x)\big] .
\end{equation}
Taking the trace over the color indices ensures that (2.17) is invariant
under gauge transformations. One can use (2.16), the invariance of the trace
under cyclic permutations, and $\Omega(x)^\dagger = \Omega(x)^{-1}$ to
verify this property. Furthermore, the summation over the Lorentz indices
$\mu$, $\nu$ (summation convention) ensures that the action is a Lorentz
scalar. We stress again that (2.17) is understood in Euclidean space.

From (2.15) and (2.17) it is obvious that our gauge action generalizes the
action of electrodynamics. Up to the trace and the different overall factor,
it exactly matches the action for the electromagnetic field. The trace is
due to the fact that gluon fields are matrix valued. Below, we decompose the
matrix-valued fields into components and in this way get rid of the trace,
pushing the similarity to electrodynamics even further. The extra factor
$1/g^2$ is just a convenient way to introduce the coupling. After rescaling
the gauge fields
\begin{equation}
  A_\mu(x) \to \frac{1}{g} A_\mu(x) ,
\end{equation}
the factor $1/g^2$ in (2.17) is gone and the gauge action assumes the more
familiar form. Now, the gauge coupling shows up in the covariant derivative
\begin{equation}
  D_\mu(x) \to \partial_\mu + \im g A_\mu(x) ,
\end{equation}
making obvious that $g$ is the coupling strength of the gauge fields to the
quarks. On the lattice it is more convenient to have the gauge coupling as
an overall factor of the gauge action, i.e., we work with the form (2.17).

\subsection{Color components of the gauge field}

We have introduced the gauge fields $A_\mu(x)$ as hermitian, traceless
matrices and have shown that the gauge transformation (2.11) preserves these
properties. Thus the $A_\mu(x)$ are in the Lie algebra su(3) and we can
write
\begin{equation}
  A_\mu(x) = \sum_{i=1}^{8} A_\mu^{(i)}(x)\, T_i .
\end{equation}
The components $A_\mu^{(i)}(x)$, $i = 1, 2, \ldots 8$, are real-valued
fields, the so-called color components, and the $T_i$ are a basis for
traceless hermitian $3 \times 3$ matrices (see Appendix A.1). We can use
this representation (2.20) of the gauge field to write also the field
strength tensor (2.15) in terms of its components. Inserting (2.20) into
(2.15) we obtain
\begin{equation}
  F_{\mu\nu}(x) = \sum_{i=1}^{8}
    \big(\partial_\mu A_\nu^{(i)}(x) - \partial_\nu A_\mu^{(i)}(x)\big) T_i
    + \im \sum_{j,k=1}^{8} A_\mu^{(j)}(x) A_\nu^{(k)}(x)\,
      [T_j, T_k] .
\end{equation}
The commutator on the right-hand side can be simplified further with the
commutation relations (A.4) and one ends up with
\begin{equation}
  F_{\mu\nu}(x) = \sum_{i=1}^{8} F_{\mu\nu}^{(i)}(x)\, T_i ,
\end{equation}
\begin{equation}
  F_{\mu\nu}^{(i)}(x) = \partial_\mu A_\nu^{(i)}(x) - \partial_\nu A_\mu^{(i)}(x)
    - f_{ijk}\, A_\mu^{(j)}(x) A_\nu^{(k)}(x) .
\end{equation}
This representation of the field strength can now be inserted in the
expression (2.17) for the gauge action and, using (A.3) to evaluate the
trace, we obtain
\begin{equation}
  S_G[A] = \frac{1}{4 g^2} \sum_{i=1}^{8}
    \int \dd^4 x\, F_{\mu\nu}^{(i)}(x) F_{\mu\nu}^{(i)}(x) .
\end{equation}
From this equation we see that the gauge action is a sum over color
components and each term has the form of the action of electrodynamics.
However, there appears a qualitatively new feature: From the right-hand side
of (2.23) we see that the field strength color components are not linear in
the gauge field $A_\mu^{(i)}(x)$ but have a quadratic piece which mixes the
different color components of the gluon field. Thus, in the action (2.24) we
not only encounter the term quadratic in the gauge fields which is familiar
from electrodynamics, but also find cubic and quartic terms. These terms
give rise to self-interactions of the gluons, making QCD a highly
nontrivial theory. The self-interactions are responsible for confinement of
color, the most prominent feature of QCD.

In Fig.~2.1 we show a schematic picture (a so-called tree-level Feynman
diagram) illustrating the cubic and quartic interaction terms. The curly
lines represent the gluons and the dots are the interaction vertices.

\begin{figure}[htbp]
  \centering
  \includegraphics[width=0.7\textwidth]{fig21.png}
  \caption*{Fig. 2.1. \emph{Schematic picture of the cubic and quartic gluon
    self-interaction. The wavy lines represent the gluon propagators and the
    dots are the interaction vertices}.}
  \label{fig:21}
\end{figure}

These remarks conclude our review of the continuum action and we have all
the concepts and notations necessary to begin the lattice discretization of
quantum chromodynamics. We stress, however, that besides gauge invariance,
there are other important symmetries of the QCD action. We discuss these
symmetries later as we need them.

\section{Naive discretization of fermions}

In this section we introduce the so-called naive discretization of the
fermion action. Later, in Chap.~5, this discretization will be augmented
with an additional term to remove lattice artifacts. Here, it serves to
present the basic idea and, more importantly, to discuss the representation
of the lattice gluon field which differs from the continuum form. We show
that on the lattice the gluon fields must be introduced as elements of the
gauge group and not as elements of the algebra, as is done in the continuum
formulation.

\subsection{Discretization of free fermions}

As discussed in Sect.~1.4, the first step in the lattice formulation is the
introduction of the 4D lattice $\Lambda$:
\begin{equation}
  \Lambda = \{n = (n_1, n_2, n_3, n_4) \;|\;
    n_1, n_2, n_3 = 0, 1, \ldots, N-1 ; \;
    n_4 = 0, 1, \ldots, N_T-1\} .
\end{equation}
The vectors $n \in \Lambda$ label points in space--time separated by a
lattice constant $a$. In our lattice discretization of QCD we now place
spinors at the lattice points only, i.e., our fermionic degrees of freedom
are
\begin{equation}
  \psi(n) , \quad \bar\psi(n) , \quad n \in \Lambda ,
\end{equation}
where the spinors carry the same color, Dirac, and flavor indices as in the
continuum (we suppress them in our notation). Note that for notational
convenience we only use the integer-valued 4-coordinate $n$ to label the
lattice position of the quarks and not the actual physical space--time point
$x = an$.

In the continuum the action $S_F^0$ for a free fermion is given by the
expression (set $A_\mu = 0$ in (2.5))
\begin{equation}
  S_F^0[\psi, \bar\psi] = \int \dd^4 x\,
    \bar\psi(x) (\gamma_\mu \partial_\mu + m) \psi(x) .
\end{equation}
When formulating this action on the lattice we have to discretize the
integral over space--time as well as the partial derivative. The
discretization is implemented as a sum over $\Lambda$, as we did for the
scalar field theory in Chap.~1. The partial derivative is discretized with
the symmetric expression
\begin{equation}
  \partial_\mu \psi(x) \to
    \frac{1}{2a}\big(\psi(n + \hat\mu) - \psi(n - \hat\mu)\big) .
\end{equation}
Thus, our lattice version of the free fermion action reads
\begin{equation}
  S_F^0[\psi, \bar\psi] = a^4 \sum_{n\in\Lambda}
    \bar\psi(n) \left( \sum_{\mu=1}^{4} \gamma_\mu\,
      \frac{\psi(n + \hat\mu) - \psi(n - \hat\mu)}{2a} + m\,\psi(n) \right) .
\end{equation}
This form is the starting point for the introduction of the gauge fields.

\subsection{Introduction of the gauge fields as link variables}

In the last section we showed that requiring the invariance of the action
under the local rotation (2.7) of the color indices of the quark fields
enforces the introduction of the gauge fields. On the lattice we implement
the same transformation by choosing an element $\Omega(n)$ of SU(3) for each
lattice site $n$ and transforming the fermion fields according to
\begin{equation}
  \psi(n) \to \psi'(n) = \Omega(n) \psi(n) , \qquad
  \bar\psi(n) \to \bar\psi'(n) = \bar\psi(n) \Omega(n)^\dagger .
\end{equation}
As for the continuum case, we find that on the lattice the mass term of
(2.29) is invariant under the transformation (2.30). For the discretized
derivative terms in (2.29), this is not the case. Consider, e.g., the term
\begin{equation}
  \bar\psi(n)\psi(n + \hat\mu) \to
  \bar\psi'(n) \psi'(n + \hat\mu) =
  \bar\psi(n) \Omega(n)^\dagger \Omega(n + \hat\mu) \psi(n + \hat\mu) .
\end{equation}
This is not gauge-invariant. If, however, we introduce a field $U_\mu(n)$
with a directional index $\mu$, then
\begin{equation}
  \bar\psi'(n) U'_\mu(n) \psi'(n + \hat\mu) =
  \bar\psi(n) \Omega(n)^\dagger U'_\mu(n) \Omega(n + \hat\mu)
  \psi(n + \hat\mu)
\end{equation}
is gauge-invariant if we define the gauge transformation of the new field by
\begin{equation}
  U_\mu(n) \to U'_\mu(n) = \Omega(n) U_\mu(n) \Omega(n + \hat\mu)^\dagger .
\end{equation}
In order to make the fermionic action (2.29) gauge-invariant, we introduce
the gauge fields $U_\mu(n)$ as elements of the gauge group SU(3) which
transform as given in (2.33). These matrix-valued variables are oriented and
are attached to the links of the lattice and thus are often referred to as
link variables. $U_\mu(n)$ lives on the link which connects the sites $n$
and $n + \hat\mu$.

Since the link variables are oriented, we can also define link variables
that point in negative $\mu$ direction. Note that these are not independent
link variables but are introduced only for notational convenience. In
particular, $U_{-\mu}(n)$ points from $n$ to $n - \hat\mu$ and is related to
the positively oriented link variable $U_\mu(n - \hat\mu)$ via the definition
\begin{equation}
  U_{-\mu}(n) \equiv U_\mu(n - \hat\mu)^\dagger .
\end{equation}
In Fig.~2.2 we illustrate the geometrical setting of the link variables on
the lattice.

\begin{figure}[htbp]
  \centering
  \includegraphics[width=0.7\textwidth]{fig22.png}
  \caption*{Fig. 2.2. \emph{The link variables $U_\mu(n)$ and $U_{-\mu}(n)$}.}
  \label{fig:22}
\end{figure}

From the definitions (2.34) and (2.33) we obtain the transformation
properties of the link in negative direction
\begin{equation}
  U_{-\mu}(n) \to U'_{-\mu}(n) =
  \Omega(n) U_{-\mu}(n) \Omega(n - \hat\mu)^\dagger .
\end{equation}
Note that we have introduced the gluon fields $U_\mu(n)$ as elements of the
gauge group SU(3), not as elements of the Lie algebra which were used in the
continuum. According to the gauge transformations (2.33) and (2.35) also the
transformed link variables are elements of the group SU(3).

Having introduced the link variables and their properties under gauge
transformations, we can now generalize the free fermion action (2.29) to the
so-called naive fermion action for fermions in an external gauge field $U$:
\begin{equation}
  \begin{split}
    S_F[\psi, \bar\psi, U] = a^4 \sum_{n\in\Lambda}
    \bar\psi(n) \Big( \sum_{\mu=1}^{4} \gamma_\mu\,
      \frac{U_\mu(n)\psi(n + \hat\mu) - U_{-\mu}(n)\psi(n - \hat\mu)}{2a}\\
    + m\,\psi(n) \Big) .
  \end{split}
\end{equation}
Using (2.30), (2.33), and (2.35) for the gauge transformation properties of
fermions and link variables, one readily sees the gauge invariance of the
fermion action (2.36), $S_F[\psi, \bar\psi, U] = S_F[\psi', \bar\psi', U']$.

\subsection{Relating the link variables to the continuum gauge fields}

Let us now discuss the link variables in more detail and see how they can be
related to the algebra-valued gauge fields of the continuum formulation. We
have introduced $U_\mu(n)$ as the link variable connecting the points $n$
and $n + \hat\mu$. The gauge transformation properties (2.33) are
consequently governed by the two transformation matrices $\Omega(n)$ and
$\Omega(n + \hat\mu)^\dagger$. Also in the continuum an object with such
transformation properties is known: It is the path-ordered exponential
integral of the gauge field $A_\mu$ along some curve $C_{xy}$ connecting two
points $x$ and $y$, the so-called gauge transporter:
\begin{equation}
  G(x, y) = \mathcal{P} \exp\!\left( \im \int_{C_{xy}} A \cdot ds \right) .
\end{equation}
We may assume that a lattice is embedded in the continuum, where smooth
gauge fields live. For a detailed discussion of the path-ordered exponential
in the continuum see, e.g., [2]. We do not need the precise definition of
the continuum gauge transporters and only use that under a gauge
transformation (2.11) they transform as
\begin{equation}
  G(x, y) \to \Omega(x) G(x, y) \Omega(y)^\dagger .
\end{equation}
These transformation properties are the same as for our link variables
$U_\mu(n)$ when $n$ and $n + \hat\mu$ are considered as end points of a
path. Based on these transformation properties, we interpret the link
variable $U_\mu(n)$ as a lattice version of the gauge transporter connecting
the points $n$ and $n + \hat\mu$, i.e., we wish to establish
$U_\mu(n) = G(n, n + \hat\mu) + O(a)$. For that purpose we introduce
algebra-valued lattice gauge fields $A_\mu(n)$ and write
\begin{equation}
  U_\mu(n) = \exp\!\big(\im a A_\mu(n)\big) .
\end{equation}
When comparing (2.37) and (2.39) one sees that we have approximated the
integral along the path from $n$ to $n + \hat\mu$ by $aA_\mu(n)$, i.e., by
the length $a$ of the path times the value of the field $A_\mu(n)$ at the
starting point.\footnote{We remind the reader that for notational
convenience we denote the lattice points only by their integer-valued
4-coordinates $n$. For the current discussion it should, however, be kept in
mind that the physical space--time coordinates are $a\,n$, i.e., neighboring
lattice points are separated by the distance $a$.} This approximation is
good to $O(a)$ and no path ordering is necessary to that order. Since the
link variables act as gauge transporters, we will often use this
nomenclature instead of ``link variable.''

Based on the relation (2.39) we can now also connect the lattice fermion
action (2.36) to its continuum counterpart (2.5). Since one of the guiding
principles of our construction is the requirement that in the limit
$a \to 0$ the lattice action approaches the continuum form, we expand (2.39)
for small $a$,
\begin{equation}
  U_\mu(n) = 1 + \im a A_\mu(n) + O(a^2) , \qquad
  U_{-\mu}(n) = 1 - \im a A_\mu(n - \hat\mu) + O(a^2) ,
\end{equation}
where we use (2.34) and $A_\mu = A_\mu^\dagger$ for the expansion of
$U_{-\mu}(n)$. Inserting these expanded link variables into expression
(2.36) we find
\begin{equation}
  S_F[\psi, \bar\psi, U] = S_F^0[\psi, \bar\psi] + S_F^I[\psi, \bar\psi, A] ,
\end{equation}
where $S_F^0$ denotes the free part of the action and the interaction part
reads
\begin{equation}
  \begin{split}
    S_F^I[\psi, \bar\psi, A] &= \im a^4 \sum_{n\in\Lambda}\sum_{\mu=1}^{4}
      \bar\psi(n) \gamma_\mu\,
      \half\big(A_\mu(n) \psi(n + \hat\mu) + A_\mu(n - \hat\mu)
        \psi(n - \hat\mu)\big) \\
    &= \im a^4 \sum_{n\in\Lambda}\sum_{\mu=1}^{4}
      \bar\psi(n) \gamma_\mu A_\mu(n) \psi(n) + O(a) .
  \end{split}
\end{equation}
In the second step we have used $\psi(n \pm \hat\mu) = \psi(n) + O(a)$ and
$A_\mu(n - \hat\mu) = A_\mu(n) + O(a)$. The two Eqs. (2.41) and (2.42)
establish that when expanding the lattice version (2.36) of the fermionic
action for $a \to 0$, we indeed recover the continuum form (2.5).

Before we continue with discretizing the gauge part of the QCD action we
stress an important conceptual point: The group-valued link variables
$U_\mu(n)$ are not merely an auxiliary construction to sneak the Lie
algebra-valued fields $A_\mu(x)$ of the continuum into the lattice
formulation. Instead, the group elements $U_\mu(n)$ are considered as the
fundamental variables which are integrated over in the path integral (see
Chap.~3). This change from algebra-valued to (compact) group-valued fields
has important consequences. In particular, the role of gauge fixing changes
considerably. We will discuss these issues in detail once we have completed
the construction of QCD on the lattice.

\section{The Wilson gauge action}

We have introduced the link variables as the basic quantities for putting
the gluon field on the lattice. Now we construct a lattice gauge action in
terms of the link variables and show that in the limit $a \to 0$ it
approaches its continuum counterpart (assuming that the lattice gauge fields
are embedded in a continuous background). This is the naive continuum limit
in contradistinction to the continuum limit of the full, integrated quantum
theory.

\subsection{Gauge-invariant objects built with link variables}

As a preparation for the construction of the gluon action let us first
discuss the transformation properties of a string of link variables along a
path of links. Let $P$ be such a path of $k$ links on the lattice connecting
the points $n_0$ and $n_1$. We define the ordered product
\begin{equation}
  P[U] = U_{\mu_0}(n_0) U_{\mu_1}(n_0 + \hat\mu_0) \ldots
    U_{\mu_{k-1}}(n_1 - \hat\mu_{k-1})
  \equiv \prod_{(n,\mu)\in P} U_\mu(n) .
\end{equation}
This object is the lattice version of the continuum gauge transporter
(2.37). Note that the path $P$ may contain link variables for both
directions $\pm\mu$.

From the transformation properties for single link variables, (2.33) and
(2.35), it follows that gauge rotations for all but the end points cancel:
Consider the transformation of two subsequent link variables on the path,
one ending at $n$ the other one starting from this point. The two
transformation matrices $\Omega(n)^\dagger$ and $\Omega(n)$ cancel each
other at $n$. Only the matrices at the two end points of the path, $n_0$
and $n_1$, remain. Thus the product $P[U]$ transforms according to
\begin{equation}
  P[U] \to P[U'] = \Omega(n_0) P[U] \Omega(n_1)^\dagger .
\end{equation}
Like for the single link term, from such a product of link variables $P[U]$
a gauge-invariant quantity can be constructed by attaching quark fields at
the starting point and at the end point,
\begin{equation}
  \bar\psi(n_0) P[U] \psi(n_1) .
\end{equation}
An alternative way of constructing a gauge-invariant product of link
variables is to choose for the path $P$ and a closed loop $L$ and to take
the trace,
\begin{equation}
  L[U] = \tr\!\left[ \prod_{(n,\mu)\in L} U_\mu(n) \right] .
\end{equation}
According to (2.44), under a gauge transformation only the matrices
$\Omega(n_0)$ and $\Omega(n_0)^\dagger$ at the point $n_0$ where the loop is
rooted remain. These matrices then cancel when taking the trace. We find
\begin{equation}
  L[U'] = \tr\!\left[ \Omega(n_0) \prod_{(n,\mu)\in L} U_\mu(n)
    \Omega(n_0)^\dagger \right] = \tr\!\left[ \prod_{(n,\mu)\in L} U_\mu(n)
    \right] = L[U] .
\end{equation}
Thus, we have established that the trace over a closed loop of link
variables is a gauge-invariant object. Such loops of link variables are used
for the construction of the gluon action and later will also serve as
physical observables.

\subsection{The gauge action}

For the gluon action it is sufficient to use the shortest, nontrivial
closed loop on the lattice, the so-called plaquette. The plaquette variable
$U_{\mu\nu}(n)$ is a product of only four link variables defined as
\begin{equation}
  \begin{split}
    U_{\mu\nu}(n) &= U_\mu(n) U_\nu(n + \hat\mu) U_{-\mu}(n + \hat\mu + \hat\nu)
      U_{-\nu}(n + \hat\nu)\\
    &= U_\mu(n) U_\nu(n + \hat\mu) U_\mu(n + \hat\nu)^\dagger U_\nu(n)^\dagger .
  \end{split}
\end{equation}
In the second formulation we have utilized the equivalence (2.34). We depict
the plaquette in Fig.~2.3. As we have shown in the last paragraph, the trace
of the plaquette variable is a gauge-invariant object.

\begin{figure}[htbp]
  \centering
  \includegraphics[width=0.7\textwidth]{fig23.png}
  \caption*{Fig. 2.3. \emph{The four link variables which build up the
    plaquette $U_{\mu\nu}(n)$. The circle indicates the order that the links
    are run through in the plaquette}.}
  \label{fig:23}
\end{figure}

We now present Wilson's form of the gauge action [3] -- the first
formulation of lattice gauge theory -- and subsequently show that it indeed
approaches the continuum form in the naive limit $a \to 0$. The Wilson
gauge action is a sum over all plaquettes, with each plaquette counted with
only one orientation. This sum can be realized by a sum over all lattice
points $n$ where the plaquettes are located, combined with a sum over the
Lorentz indices $1 \leq \mu < \nu \leq 4$,
\begin{equation}
  S_G[U] = \frac{2}{g^2} \sum_{n\in\Lambda} \sum_{\mu<\nu}
    \Re\, \tr\big[1 - U_{\mu\nu}(n)\big] .
\end{equation}
The individual contributions are the real parts of traces over the unit
matrix minus the plaquette variable. The factor $2/g^2$ is included to match
the form of the continuum action (2.17) in the limit $a \to 0$. Let us now
discuss this limit.

For establishing the correct limit we need to expand the link variables in
the form (2.39) for small $a$. In order to handle the products of the four
link variables in the plaquette in an organized way, it is useful to invoke
the Baker--Campbell--Hausdorff formula for the product of exponentials of
matrices:
\begin{equation}
  \exp(A) \exp(B) = \exp\!\left( A + B + \half [A, B] + \cdots \right) ,
\end{equation}
where $A$ and $B$ are arbitrary matrices and the dots indicate powers of the
matrices larger than 2 which are omitted. The formula (2.50) can be proven
easily by expanding both sides in powers of $A$ and $B$. Inserting (2.39)
into the definition (2.48) of the plaquette and applying (2.50) iteratively,
we obtain
\begin{multline}
  U_{\mu\nu}(n) = \exp\!\Big(
    \im a A_\mu(n) + \im a A_\nu(n + \hat\mu)
    - \frac{a^2}{2} [A_\mu(n), A_\nu(n + \hat\mu)] \\
    - \im a A_\mu(n + \hat\nu) - \im a A_\nu(n)
    - \frac{a^2}{2} [A_\mu(n + \hat\nu), A_\nu(n)] \\
    + \frac{a^2}{2} [A_\nu(n + \hat\mu), A_\mu(n + \hat\nu)]
    + \frac{a^2}{2} [A_\mu(n), A_\nu(n)] \\
    + \frac{a^2}{2} [A_\mu(n), A_\mu(n + \hat\nu)]
    + \frac{a^2}{2} [A_\nu(n + \hat\mu), A_\nu(n)] + O(a^3) \Big) .
\end{multline}
In this expression we have gauge fields with shifted arguments such as
$A_\nu(n + \hat\mu)$. We now perform a Taylor expansion for these fields,
i.e., we set
\begin{equation}
  A_\nu(n + \hat\mu) = A_\nu(n) + a\, \partial_\mu A_\nu(n) + O(a^2) ,
\end{equation}
in all these terms and take into account contributions up to $O(a^2)$. With
this expansion several terms cancel and we obtain
\begin{equation}
  \begin{split}
    U_{\mu\nu}(n) &= \exp\!\big( \im a^2
      \big(\partial_\mu A_\nu(n) - \partial_\nu A_\mu(n)
      + \im [A_\mu(n), A_\nu(n)]\big) + O(a^3) \big) \\
    &= \exp\!\big( \im a^2 F_{\mu\nu}(n) + O(a^3) \big) .
  \end{split}
\end{equation}
In the second step we use the continuum definition of the field strength
given in (2.15). The form (2.53) can now be inserted in (2.49) for the
Wilson gauge action. The exponential in (2.53) is expanded and we find
\begin{equation}
  \begin{split}
    S_G[U] &= \frac{2}{g^2} \sum_{n\in\Lambda} \sum_{\mu<\nu}
      \Re\, \tr\big[1 - U_{\mu\nu}(n)\big] \\
    &= \frac{a^4}{2 g^2} \sum_{n\in\Lambda} \sum_{\mu,\nu}
      \tr\big[F_{\mu\nu}(n)^2\big] + O(a^2) .
  \end{split}
\end{equation}
The terms of $O(a^2)$ that appear in the expansion of the exponential in
(2.53) cancel when taking the real part of $\tr[1 - U_{\mu\nu}(n)]$ (one may
use $\tr[U_{\mu\nu}(n)]^* = \tr[U_{\mu\nu}(n)^\dagger] = \tr[U_{\nu\mu}(n)]$
to see this). In a similar way also the $O(a^3)$ terms in the expansion of
(2.53) cancel, such that the Wilson action approximates the continuum form
up to $O(a^2)$, as stated in (2.54). Note that the factor $a^4$ together
with the sum over $\Lambda$ is just the discretization of the space--time
integral and thus $\lim_{a\to 0} S_G[U] = S_G[A]$. This completes our
discussion of the naive continuum limit $a \to 0$ for the Wilson gauge
action.

Concluding this section we remark that different lattice actions for the
gauge fields have been proposed in order to reduce cutoff effects further.
We return to this issue in Chap.~9.

\section{Formal expression for the QCD lattice path integral}

Having constructed the fermion and gauge field actions on the lattice we can
now write down the complete expression for the lattice QCD path integral
formula for Euclidean correlators. As already announced, the preliminary
formulation presented here will be refined further in subsequent chapters.
Nevertheless, we find such an intermediate summary helpful for understanding
the line of the construction.

\subsection{The QCD lattice path integral}

Following the discussion of Sect.~1.4 we write Euclidean correlators as a
lattice path integral in the form
\begin{equation}
  \ev{O_2(t)\, O_1(0)} = \frac{1}{Z} \int
    \mathcal{D}[\psi, \bar\psi]\, \mathcal{D}[U]\,
    e^{-S_F[\psi,\bar\psi,U] - S_G[U]}\,
    O_2[\psi, \bar\psi, U]\, O_1[\psi, \bar\psi, U] ,
\end{equation}
where the partition function $Z$ is given by
\begin{equation}
  Z = \int \mathcal{D}[\psi, \bar\psi]\, \mathcal{D}[U]\,
    e^{-S_F[\psi,\bar\psi,U] - S_G[U]} .
\end{equation}
The quantization of the system in the path integral formalism is implemented
as an integral over all field configurations. On the lattice the
corresponding path integral measures are products of measures of all quark
field components and products of measures for all link variables:
\begin{equation}
  \begin{split}
    \mathcal{D}[\psi, \bar\psi] &=
      \prod_{n\in\Lambda}\prod_{f,\alpha,c}
      \dd\psi^{(f)}(n)_c^\alpha\,
      \dd\bar\psi^{(f)}(n)_c^\alpha ,\\
    \mathcal{D}[U] &= \prod_{n\in\Lambda}\prod_{\mu=1}^{4} \dd U_\mu(n) .
  \end{split}
\end{equation}
Both the fermion and gauge field measures shown here will be discussed in
more detail in subsequent sections. For the fermion fields we have to
include the Pauli's principle, turning the spinors $\psi$ and $\bar\psi$
into anticommuting variables. These so-called Grassmann numbers and the
corresponding rules of integration will be discussed in Chap.~5. For the
gauge fields we have denoted in (2.57) the measure for a single link
variable as $\dd U_\mu(n)$, but not yet specified how we implement the
integration over the group manifold of SU(3). This leads to the concept of
Haar measure which we discuss in the next chapter. Despite these issues, the
expressions in (2.57) already incorporate some of the essential features of
lattice QCD, in particular the reduction of the original quantum fields to a
countable number of classical variables.

As discussed in Sect.~1.4, in the path integral quantization (2.55) the
configurations to be integrated over are weighted with the Boltzmann factor
of the Euclidean action. The corresponding lattice versions of the fermion
and gauge parts of the action have been derived as (compare (2.36) and
(2.49))
\begin{multline}
  S_F[\psi, \bar\psi, U] = a^4 \sum_{f=1}^{\Nf}\sum_{n\in\Lambda}
    \bar\psi^{(f)}(n) \Big( \sum_{\mu=1}^{4} \gamma_\mu\\
    \frac{U_\mu(n)\psi^{(f)}(n + \hat\mu)
      - U_{-\mu}(n)\psi^{(f)}(n - \hat\mu)}{2a}
    + m^{(f)}\, \psi^{(f)}(n) \Big) \\
    + \text{terms discussed in Chap.~5} ,
\end{multline}
where we now also sum over $\Nf$ flavors of quarks. We stress again that the
fermion action has to be augmented with another term in order to remove
lattice artifacts (see Chap.~5). The gauge action, however, is ready to go
and is taken over unchanged from (2.49):
\begin{equation}
  S_G[U] = \frac{2}{g^2} \sum_{n\in\Lambda} \sum_{\mu<\nu}
    \Re\, \tr\big[1 - U_{\mu\nu}(n)\big] .
\end{equation}
Let us remark that the functionals $O_1[\psi, \bar\psi, U]$ and
$O_2[\psi, \bar\psi, U]$ are translations of the operators $\hat{O}_1$ and
$\hat{O}_2$ acting in Hilbert space. The translation proceeds as expressed
in (1.81) by evaluating the operators between field eigenstates. We stress
that the functional $O_2$ depends only on the fields with time argument
$n_t$ related to Euclidean time $t$ on the right-hand side of (2.55) via
$t = a\, n_t$. The fields in the functional $O_1$ depend only on the fields
with $t = 0$.

Equations (2.55), (2.56), (2.57), (2.58), and (2.59) comprise our current
status in the construction of lattice QCD. So far we discussed the
fundamental fields describing quarks and gluons and put them onto the
lattice -- quarks on the sites and the gauge fields on the links of the
lattice. In order to allow for color rotations of the fermion fields in the
same way as in the continuum, the gauge fields were introduced as
group-valued link variables. Subsequently, we showed that products of link
variables along closed paths are gauge-invariant and we have used the
plaquette to construct Wilson's gauge action. Finally, we put together these
ingredients according to the path integral quantization recipe of Sect.~1.4.
Operators are translated into functionals of classical fields. These
functionals are weighted with the Boltzmann factor and this product is then
integrated over all possible field configurations. The precise definition of
this integration -- Grassmann integration for the fermions, Haar measure for
the link variables -- will be discussed in subsequent chapters.

\section*{References}

\begin{enumerate}[label={[\arabic*]}]
  \item C.~N.~Yang and R.~Mills: Phys.\ Rev.\ 96, 191 (1954).
  \item M.~E.~Peskin and D.~V.~Schroeder: \emph{An Introduction to Quantum
        Field Theory} (Addison-Wesley, Reading, Massachusetts 1995).
  \item K.~G.~Wilson: Phys.\ Rev.\ D 10, 2445 (1974).
\end{enumerate}
